\section*{第 \ref{ch:b2:fourier} 章 --- 傅里叶级数}

\begin{solution}{exo:b2:fourier:1}
奇性使 $a_n$ 全为零。对 $n \geq 1$：
\[
b_n = \frac{2}{\pi}\int_0^{\pi} \sin nt\,\dd t
= \frac{2}{\pi}\cdot\frac{1 - (-1)^n}{n}
= \begin{cases} \frac{4}{\pi n} & n \text{ 奇数},\\ 0 & n
\text{ 偶数}. \end{cases}
\]
于是 $S(f)(t) = \frac{4}{\pi}\sum_{k\geq0} \frac{\sin\bigl((2k+1)t
\bigr)}{2k+1}$。由狄利克雷定理，在 $t = \frac\pi2$ 处（连续点，
取值 $1$）：$\sin\bigl((2k+1)\frac\pi2\bigr) = (-1)^k$，得
\[
1 = \frac4\pi \sum_{k\geq0}\frac{(-1)^k}{2k+1}
\quad\Longrightarrow\quad
\sum_{k\geq0}\frac{(-1)^k}{2k+1} = \frac{\pi}{4}
\quad\text{（莱布尼茨）}.
\]
在跳跃点 $t = 0$ 处，级数之和为 $0 = \frac{f(0^+) +
f(0^-)}{2}$，与狄利克雷的规定相符（各项均为零，
因而相容）。
\end{solution}

\begin{solution}{exo:b2:fourier:2}
偶性使 $b_n$ 全为零；$a_0 = \frac{1}{\pi}\int_{-\pi}^\pi\abs
t\,\dd t = \pi$，对 $n \geq 1$：
\[
a_n = \frac{2}{\pi}\int_0^\pi t\cos nt\,\dd t
= \frac{2}{\pi}\cdot\frac{(-1)^n - 1}{n^2}
= \begin{cases} -\frac{4}{\pi n^2} & n \text{ 奇数},\\ 0 & n
\text{ 偶数}, \end{cases}
\]
（一次分部积分）。于是
\[
\abs t = \frac{\pi}{2} - \frac{4}{\pi}\sum_{k\geq0}
\frac{\cos\bigl((2k+1)t\bigr)}{(2k+1)^2} ,
\]
级数\emph{正规收敛}（$\sum (2k+1)^{-2} < \infty$）---
正如\cref{thm:b2:fourier:parseval}~(1) 对此连续
分段 $C^1$ 函数所预言的。当 $t = 0$ 时：
\[
0 = \frac\pi2 - \frac4\pi\sum_{k\geq0}\frac{1}{(2k+1)^2}
\quad\Longrightarrow\quad
\sum_{k\geq0}\frac{1}{(2k+1)^2} = \frac{\pi^2}{8} .
\]
将 $\sum \frac{1}{n^2}$ 按奇偶项拆分：$S =
\frac{\pi^2}{8} + \frac S4$，故 $S = \frac{\pi^2}{6}$，这又是巴塞尔问题的结果。
\end{solution}

\begin{solution}{exo:b2:fourier:3}
偶性：$b_n = 0$；$a_0 = \frac{1}{\pi}\int_{-\pi}^{\pi} t^2 =
\frac{2\pi^2}{3}$；两次分部积分得 $a_n =
\frac{4(-1)^n}{n^2}$（$n \geq 1$）。于是
\[
t^2 = \frac{\pi^2}{3} + 4\sum_{n\geq1}
\frac{(-1)^n}{n^2}\cos nt
\qquad (\abs t \leq \pi),
\]
此级数正规收敛。当 $t = 0$ 时：$0 = \frac{\pi^2}{3} +
4\sum\frac{(-1)^n}{n^2}$，即 $\sum_{n\geq1}
\frac{(-1)^{n+1}}{n^2} = \frac{\pi^2}{12}$。由帕塞瓦尔等式：
\[
\frac{1}{2\pi}\int_{-\pi}^{\pi} t^4\,\dd t = \frac{\pi^4}{5}
= \frac{a_0^2}{4} + \frac12\sum_{n\geq1} a_n^2
= \frac{\pi^4}{9} + 8\sum_{n\geq1}\frac{1}{n^4} ,
\]
故 $\sum \frac{1}{n^4} = \frac18\bigl(\frac{\pi^4}{5} -
\frac{\pi^4}{9}\bigr) = \frac{\pi^4}{90}$。
\end{solution}

\begin{solution}{exo:b2:fourier:4}
若 $f$ 还是分段 $C^1$ 函数，则由帕塞瓦尔等式（已证）得
$\norm f_2^2 = \sum\abs{c_n}^2 = 0$，连续函数 $\abs f^2$ 的积分的严格正性
迫使 $f = 0$。

对仅为连续的 $f$，在本阶段不加证明地承认一般的帕塞瓦尔等式：
其证明同样只需一行。（稠密性路线概要：费耶尔、魏尔斯特拉斯型
三角逼近表明三角多项式
在连续周期函数中是 $\norm\cdot_2$-稠密的；由于 $f
\perp$ 与所有这些函数正交，$\norm f_2^2 = \langle f, f - P\rangle \leq
\norm f_2\norm{f - P}_2$ 对逼近函数 $P$ 成立，迫使 $\norm f_2 =
0$。）
\end{solution}

\begin{solution}{exo:b2:fourier:5}
偶性：$b_n = 0$；
\[
a_n = \frac{2}{\pi}\int_0^\pi \cos(\alpha t)\cos(nt)\,\dd t
= \frac{2}{\pi}\cdot
\frac{(-1)^n\,\alpha\sin(\pi\alpha)}{\alpha^2 - n^2}
\]
（先积化和差，再积分；$a_0 =
\frac{2\sin(\pi\alpha)}{\pi\alpha}$）。在 $t = \pi$ 处（
周期延拓的一个连续点，其单侧取值
由偶性而相同）：
\[
\cos(\pi\alpha)
= \frac{\sin(\pi\alpha)}{\pi\alpha}
+ \sum_{n\geq1} \frac{2\alpha\sin(\pi\alpha)}{\pi(\alpha^2 -
n^2)}\,(-1)^n\cos(n\pi)
= \frac{\sin(\pi\alpha)}{\pi}\Bigl(\frac{1}{\alpha} +
\sum_{n\geq1}\frac{2\alpha}{\alpha^2 - n^2}\Bigr),
\]
利用 $(-1)^n\cos n\pi = 1$。两边除以 $\sin(\pi\alpha)/\pi$
即得余切展开式。
\end{solution}

\begin{solution}{exo:b2:fourier:6}
反复使用 $c_n(f') = \iu n\,c_n(f)$（$k$ 次，
在 $C^{k}$ 各段上分部积分，边界值匹配）：
$c_n(f^{(k)}) = (\iu n)^k c_n(f)$。分段连续函数 $f^{(k)}$ 的系数
是有界的（事实上 $\to 0$，由贝塞尔不等式）：
\[
\abs{c_n(f)} = \frac{\abs{c_n(f^{(k)})}}{\abs n^k}
= O\bigl(\abs n^{-k}\bigr) .
\]
\end{solution}

\begin{solution}{exo:b2:fourier:7}
对 $f$ 与 $f'$ 分别使用帕塞瓦尔恒等式（两者均适用：$f$ 为 $C^1$，$f'$
分段连续 --- 事实上连续）：
\[
\frac{1}{2\pi}\int \abs f^2 = \sum_{n\neq0} \abs{c_n}^2
\quad (c_0 = 0 \text{ 由零均值假设}),
\qquad
\frac{1}{2\pi}\int \abs{f'}^2 = \sum_{n\neq0} n^2\abs{c_n}^2 .
\]
逐项比较，$n^2\abs{c_n}^2 \geq \abs{c_n}^2$ 对 $\abs n \geq 1$ 成立：
由此得出不等式。等号成立时必有 $(n^2 - 1)\abs{c_n}^2 = 0$
对所有 $n$，即 $c_n = 0$（$\abs n \geq 2$）：$f(t) =
c_1\eu^{\iu t} + c_{-1}\eu^{-\iu t} = a\cos t + b\sin t$（实
形式）；反之，满足此形式的 $f$ 使等号成立。
\end{solution}

\begin{solution}{exo:b2:fourier:8}
由 \cref{exo:b2:fourier:1} 可知 $S_{2N-1}(x) =
\frac4\pi\sum_{k=0}^{N-1}\frac{\sin((2k+1)x)}{2k+1}$。取 $x_N =
\frac{\pi}{2N}$，记 $u_k = (2k+1)x_N = \frac{(2k+1)\pi}{2N}$：
\[
S_{2N-1}(x_N) = \frac{4}{\pi}\sum_{k=0}^{N-1}\frac{\sin u_k}{2k+1}
= \frac{4}{\pi}\sum_{k=0}^{N-1}\frac{\sin u_k}{u_k}\cdot
\frac{u_k}{2k+1}
= \frac{2}{\pi}\sum_{k=0}^{N-1}\frac{\sin u_k}{u_k}\cdot
\frac{\pi}{N},
\]
由 $\frac{u_k}{2k+1} = \frac{\pi}{2N}$ 得，点 $u_k$ 恰是 $N$
等分 $\intcc{0}{\pi}$ 所得子区间的中点，每段长
$\frac{\pi}{N}$，该和为
连续函数 $u \mapsto \frac{\sin u}{u}$ 的中点黎曼和，因此收敛于
\[
\frac{2}{\pi}\int_0^\pi \frac{\sin u}{u}\,\dd u
\approx \frac{2}{\pi}\times 1.8519 \approx 1.179 .
\]
跳跃点附近的部分和超出值 $1$ 约 $\approx
18\%$ 的半跳跃量，并且始终如此：吉布斯现象的定量刻画。
\end{solution}

\begin{solution}{exo:b2:fourier:9}
由 $\cos 3t = 4\cos^3t - 3\cos t$ 得：
\[
\cos^3 t = \frac{3\cos t + \cos 3t}{4},
\qquad
\sin^2t\,\cos t = \cos t - \cos^3 t
= \frac{\cos t - \cos 3t}{4} .
\]
二者都是三角多项式，因此各自等于它自己的
傅里叶级数（系数的唯一性：两种展开式
之差将是一个
系数全为零的三角多项式）。对 $\cos^3t$：$a_1 = \frac34$，$a_3 =
\frac14$，其余 $a_n$ 及全部 $b_n$ 为零；$c_{\pm1} =
\frac38$，$c_{\pm3} = \frac18$。对 $\sin^2t\cos t$：$a_1 =
\frac14$，$a_3 = -\frac14$；$c_{\pm1} = \frac18$，$c_{\pm3} =
-\frac18$。
\end{solution}

\begin{solution}{exo:b2:fourier:10}
直接计算：
\[
c_n = \frac{1}{2\pi}\int_{-\pi}^{\pi}\eu^{(a - \iu n)t}\dd t
= \frac{\eu^{(a-\iu n)\pi} - \eu^{-(a - \iu n)\pi}}
{2\pi(a - \iu n)}
= \frac{(-1)^n\sinh(a\pi)}{\pi(a - \iu n)} ,
\]
利用 $\eu^{\pm\iu n\pi} = (-1)^n$。在 $t = \pi$ 处，周期延拓
由 $\eu^{a\pi}$ 跳至 $\eu^{-a\pi}$；狄利克雷定理
（取对称部分和）给出
\[
\cosh(a\pi) = \sum_{n\in\Z}(-1)^n c_n\,
= \frac{\sinh(a\pi)}{\pi}\Bigl(\frac1a +
\sum_{n\geq1}\Bigl(\frac{1}{a - \iu n} +
\frac{1}{a + \iu n}\Bigr)\Bigr)
= \frac{\sinh(a\pi)}{\pi}\Bigl(\frac1a +
\sum_{n\geq1}\frac{2a}{a^2 + n^2}\Bigr) ,
\]
成对各项的虚部相互抵消，两边除以
$\sinh(a\pi)$：
\[
\coth(\pi a) = \frac{1}{\pi a} +
\sum_{n\geq1}\frac{2a}{\pi(a^2 + n^2)} ,
\]
这是 \cref{exo:b2:fourier:5} 的双曲版本。
\end{solution}

\begin{solution}{exo:b2:fourier:11}
交换律：作代换 $s = x - t$，并利用被积函数的
周期性。对 $c_n(f * g)$，被积函数 $(x, t) \mapsto
f(x-t)g(t)\eu^{-\iu nx}$ 连续；两个累次积分
相等（二者均视为外层变量上限的函数，
在左端点取零，并有相同的
导数 --- 由连续性加上
\cref{thm:b2:integration:continuity} 可保证对累次积分求导
是合理的）。因此
\[
c_n(f*g) = \frac{1}{2\pi}\int_{-\pi}^{\pi} g(t)\,\eu^{-\iu nt}
\Bigl(\frac{1}{2\pi}\int_{-\pi}^{\pi}
f(x-t)\,\eu^{-\iu n(x-t)}\dd x\Bigr)\dd t
= c_n(f)\,c_n(g),
\]
内层积分为 $c_n(f)$，对每个 $t$ 均成立（代换与
周期性）。最后由 \cref{lem:b2:fourier:kernel} 可知
$S_N(f)(x) = \frac{1}{2\pi}\int f(x+u)D_N(u)\dd u$；
作代换 $u \mapsto -t$，并利用 $D_N$ 的偶性，可将其化为
$(f * D_N)(x)$。
\end{solution}

\begin{solution}{exo:b2:fourier:12}
\begin{enumerate}
  \item 令 $w = r\eu^{\iu t}$：
        \[
        P_r(t) = 1 + 2\,\Re\frac{w}{1 - w}
        = \Re\frac{1 + w}{1 - w}
        = \frac{1 - \abs w^2}{\abs{1 - w}^2}
        = \frac{1 - r^2}{1 - 2r\cos t + r^2} > 0 .
        \]
        均值为 $1$：对正规收敛的
        级数逐项积分，仅 $n = 0$ 项保留。
  \item 当 $\delta \leq \abs t \leq \pi$ 时，$\cos t \leq
        \cos\delta$，故 $P_r(t) \leq \frac{1 - r^2}{1 -
        2r\cos\delta + r^2}$，其分母趋于 $2 -
        2\cos\delta > 0$，而分子趋于 $0$：
        一致收敛于 $0$，即在
        $0$ 的任一邻域之外。
  \item 逐项积分（对 $t$ 的正规收敛）
        得 $(f * P_r)(x) = \sum_n r^{\abs
        n}c_n(f)\eu^{\iu nx}$。逼近单位元
        论证：利用均值为 $1$ 及非负性，
        \[
        \abs{(f*P_r)(x) - f(x)}
        \leq \frac{1}{2\pi}\int_{-\pi}^{\pi}
        \abs{f(x-t) - f(x)}\,P_r(t)\,\dd t ,
        \]
        在 $\abs t = \delta$ 处分段：至多为 $\varepsilon$
        （由海涅定理）加上 $2\norm f_\infty\sup_{\delta\leq\abs
        t\leq\pi}P_r \to \varepsilon$，故当
        $r \to 1^-$ 时一致收敛。这是傅里叶级数的阿贝尔求和，
        即幂级数一章中边界理论的
        傅里叶版本。
\end{enumerate}
\end{solution}

\begin{solution}{pb:b2:fourier:1}
\textbf{1.} 由\cref{lem:b2:fourier:kernel}，对 $n = 0,
\dots, N-1$ 取平均（利用积分的线性性）得 $\sigma_N(f)(x) =
\frac{1}{2\pi}\int f(x+u)F_N(u)\dd u$；各 $D_n$ 的均值
为 $1$，故 $F_N$ 的均值为 $1$。

\textbf{2.} 利用 $\eu^{\iu u} - 1 = 2\iu\,\eu^{\iu
u/2}\sin\frac u2$，有：
\[
\sum_{n=0}^{N-1}\sin\Bigl(\Bigl(n + \frac12\Bigr)u\Bigr)
= \Im\Bigl[\eu^{\iu u/2}\,\frac{\eu^{\iu Nu} - 1}{\eu^{\iu u}
- 1}\Bigr]
= \Re\,\frac{1 - \eu^{\iu Nu}}{2\sin\frac u2}
= \frac{1 - \cos Nu}{2\sin\frac u2}
= \frac{\sin^2\frac{Nu}2}{\sin\frac u2} .
\]
除以 $N\sin\frac u2$，得：
\[
F_N(u) = \frac1N\sum_{n=0}^{N-1}
\frac{\sin\bigl((n+\frac12)u\bigr)}{\sin\frac u2}
= \frac{1}{N}\,
\frac{\sin^2\frac{Nu}{2}}{\sin^2\frac u2} \geq 0 .
\]

\textbf{3.} 当 $\delta \leq \abs u \leq \pi$ 时：$\sin^2\frac u2
\geq \sin^2\frac\delta2$，且 $\sin^2\frac{Nu}2 \leq 1$：$F_N
\leq \frac{1}{N\sin^2(\delta/2)} \to 0$，且在该区域上一致。

\textbf{4.} 由单位均值可得 $\sigma_Nf(x) - f(x) =
\frac{1}{2\pi}\int\bigl(f(x+u) - f(x)\bigr)F_N(u)\dd u$。给定
$\varepsilon$，由海涅定理，存在 $\delta$，使得 $\abs{f(x+u) - f(x)}
\leq \varepsilon$ 对 $\abs u \leq \delta$ 成立，且关于 $x$ 一致。
于是，利用 $F_N \geq 0$ 及其单位均值，
\[
\abs{\sigma_Nf(x) - f(x)}
\leq \varepsilon +
2\norm f_\infty\cdot\frac{1}{2\pi}
\int_{\delta\leq\abs u\leq\pi}F_N
\leq \varepsilon + \frac{2\norm
f_\infty}{N\sin^2\frac\delta2}
\leq 2\varepsilon
\]
对充分大的 $N$ 成立，且关于 $x$ 一致：此即费耶尔定理。

\textbf{5.} $F_N$ 为偶函数，均值为 $1$：每一半
$\intcc{0}{\pi}$、$\intcc{-\pi}{0}$ 各自贡献的均值为 $\frac12$。
于是
\[
\sigma_Nf(x) - \frac{f(x^+)+f(x^-)}{2}
= \frac{1}{2\pi}\int_0^\pi\bigl(f(x+u) -
f(x^+)\bigr)F_N\,\dd u
+ \frac{1}{2\pi}\int_{-\pi}^0\bigl(f(x+u) -
f(x^-)\bigr)F_N\,\dd u ;
\]
在每一半上于 $\abs u = \delta$ 处分段，此处单侧
极限的误差在 $\varepsilon$ 之内，并由第 3 题使远端
部分消失：两个积分均趋于 $0$。估计：由 $F_N \geq 0$
可得 $\abs{\sigma_Nf(x)} \leq
\frac{1}{2\pi}\int\abs{f(x+u)}F_N \leq \norm f_\infty$。

\textbf{6.} 每个 $\sigma_N(f)$ 都是三角多项式
（即 $S_n(f)$，$n < N$ 的平均），且 $\sigma_N(f) \to f$
一致收敛：此即三角形式的魏尔斯特拉斯定理。

\textbf{7.} 若 $c_n(f) = 0$ 对一切 $n$ 成立，则每个 $S_n(f) = 0$，
从而每个 $\sigma_N(f) = 0$；由费耶尔定理，$f = \lim\sigma_Nf =
0$。将此结论用于两函数之差：连续周期
函数由其傅里叶系数唯一确定。

\textbf{8.} $S_Nf$ 是 $f$ 到
$\mathcal T_N$ 上的正交投影（见\cref{prop:b2:fourier:bessel}），因此它使
$\norm{f - P}_2$ 在 $P \in \mathcal T_N$ 上取最小值；由于 $\sigma_Nf
\in \mathcal T_{N-1} \subseteq \mathcal T_N$，于是
\[
\norm{f - S_Nf}_2 \leq \norm{f - \sigma_Nf}_2
\leq \norm{f - \sigma_Nf}_\infty \to 0
\]
（中间的不等式成立，因为 $\abs\cdot^2$ 的均值
不超过其上确界的平方）。由勾股定理，$\norm f_2^2 = \sum_{\abs
n\leq N}\abs{c_n}^2 + \norm{f - S_Nf}_2^2$，取极限后
即得对每个连续周期函数 $f$ 的帕塞瓦尔等式。

\textbf{9.} 设 $t_1, \dots, t_p$ 为 $f$ 在一个
周期内的跳跃点，$M = \norm f_\infty$。对小的 $\eta$，定义 $g = f$
在区间 $\intoo{t_j - \eta}{t_j + \eta}$ 之外，并在每个此类区间上以
仿射弦代替：$g$ 连续，
周期，$\norm g_\infty \leq M$，且
\[
\norm{f - g}_2^2 \leq \frac{1}{2\pi}\,p\cdot(2M)^2\cdot2\eta
\leq \varepsilon^2
\]
对充分小的 $\eta$ 成立。由贝塞尔不等式，$S_N$ 在
$\norm\cdot_2$ 下为压缩映射，于是
\[
\norm{f - S_Nf}_2
\leq \norm{f - g}_2 + \norm{g - S_Ng}_2 + \norm{S_N(g -
f)}_2
\leq 2\varepsilon + \norm{g - S_Ng}_2 ,
\]
再由第 8 题得 $\limsup_N\norm{f - S_Nf}_2 \leq
2\varepsilon$ 对每个 $\varepsilon$ 成立：$\norm{f - S_Nf}_2 \to
0$，并由勾股定理，对每个分段
连续的 $f$ 成立帕塞瓦尔等式。此前在
\cref{thm:b2:fourier:parseval} 中“承认”的结论现已成为定理。

\textbf{10.} 方波满足 $\norm f_\infty = 1$，故
由第 5 题，$\abs{\sigma_N(f)} \leq 1$ 处处成立，且对
每个 $N$ 都成立 --- 绝无过冲 --- 而
由\cref{exo:b2:fourier:8}，$\sup_xS_{2N-1}(f)(x) \to
\approx 1.179$。起作用的性质是：$F_N \geq 0$，因此
$\sigma_Nf(x)$ 是一个加权\emph{平均}，其对象是 $f$ 的取值，
故永不会超出 $\intcc{\min f}{\max f}$；而 $D_N$ 取
负值，因此 $S_N$ 可以。

\textbf{11.} $\abs{\sin N\theta} \leq N\abs{\sin\theta}$ 由
归纳法（$\abs{\sin(N{+}1)\theta} \leq
\abs{\sin N\theta}\abs{\cos\theta} +
\abs{\cos N\theta}\abs{\sin\theta} \leq
(N+1)\abs{\sin\theta}$）：令 $\theta = \frac u2$，
\[
F_N(u) = \frac{\sin^2\frac{Nu}2}{N\sin^2\frac u2}
\leq \frac{N^2\sin^2\frac u2}{N\sin^2\frac u2} = N .
\]
由 $\sin$ 在 $\intcc{0}{\frac\pi2}$ 上的凹性可得
$\sin\frac u2 \geq \frac{u}{\pi}$，当 $0 \leq u \leq \pi$ 时，于是
$F_N(u) \leq \frac{1}{N(u/\pi)^2} = \frac{\pi^2}{Nu^2}$.

\textbf{12.} 由偶性及在 $\frac1N$ 处分段：
\[
\frac{1}{2\pi}\int_{-\pi}^{\pi}\abs uF_N
= \frac1\pi\int_0^\pi uF_N
\leq \frac1\pi\Bigl(\int_0^{1/N}uN\,\dd u +
\int_{1/N}^{\pi}\frac{\pi^2}{Nu}\,\dd u\Bigr)
= \frac{1}{2\pi N} + \frac{\pi\ln(\pi N)}{N} .
\]
当 $N \geq 2$ 时，$\ln(\pi N) \leq \bigl(1 +
\frac{\ln\pi}{\ln2}\bigr)\ln N \leq 2.66\ln N$，且
$\frac{1}{2\pi N} \leq \frac{\ln N}{N}$，故该矩满足
$\leq \frac{9\ln N}{N}$.

\textbf{13.} 对 $L$-利普希茨函数 $f$：
\[
\abs{\sigma_Nf(x) - f(x)}
\leq \frac{1}{2\pi}\int\abs{f(x+u) - f(x)}F_N(u)\dd u
\leq L\cdot\frac{1}{2\pi}\int\abs uF_N
\leq \frac{9L\ln N}{N},
\]
关于 $x$ 一致成立。

\textbf{14.} 由 $c_0(e_1) = 0$ 得 $S_0(e_1) = 0$，而
$S_n(e_1) = e_1$，当 $n \geq 1$ 时，$\sigma_N(e_1) =
\frac{N-1}{N}e_1$，于是 $\norm{\sigma_Ne_1 - e_1}_\infty =
\frac1N$。即使对这一整函数、带限信号，速率仍为
$\frac1N$。费耶尔方法饱和，恰如伯恩斯坦算子
在 $\frac1n$ 处饱和（沃罗诺夫斯卡娅，见函数列一章的
周末问题）。

\textbf{15.} 若 $f$ 在 $\intoo{x-\delta}{x+\delta}$ 上恒为零，
则 $\sigma_Nf(x) = \frac{1}{2\pi}\int_{\delta \leq \abs u
\leq \pi}f(x+u)F_N(u)\dd u$，其绝对值至多为
$\frac{\norm f_\infty}{N\sin^2(\delta/2)} \to 0$：切萨罗
均值在 $x$ 处只依赖 $f$ 在 $x$ 附近的取值。

\textbf{16.} 给定 $\intcc ab$ 与 $\varepsilon$，取
连续的 $1$ 周期分段仿射函数 $\varphi^\pm$，满足
$\varphi^- \leq \mathbf 1_{\intcc ab} \leq \varphi^+$，且
$\int_0^1(\varphi^+ - \varphi^-) \leq \varepsilon$（梯形，
斜率在总长为 $\varepsilon$ 的区间上）。
于是
\[
\limsup_N\frac{\#\{n \leq N : x_n \in \intcc ab\}}{N}
\leq \lim_N\frac1N\sum_{n\leq N}\varphi^+(x_n)
= \int_0^1\varphi^+ \leq b - a + \varepsilon ,
\]
同理可得 $\liminf \geq b - a - \varepsilon$，
于是这一比例趋于 $b - a$。

\textbf{17.} 对 $f = \eu^{2\iu\pi k\cdot}$，当 $k \neq 0$ 时
由假设可得极限 $0 = \int_0^1f$；当 $k = 0$ 时两边
都为 $1$；线性性处理每个三角
多项式。对连续的 $1$ 周期函数 $f$ 及 $\varepsilon >
0$，由问题 6（经 $t = 2\pi x$ 变换）得
三角多项式 $P$，使 $\norm{f - P}_\infty \leq
\varepsilon$：
\[
\Bigl|\frac1N\sum_{n\leq N}f(x_n) - \int_0^1f\Bigr|
\leq 2\varepsilon +
\Bigl|\frac1N\sum_{n\leq N}P(x_n) - \int_0^1P\Bigr|
\longrightarrow 2\varepsilon .
\]
结合问题 16 可知 $(x_n)$ 是等分布的。

\textbf{18.} 当 $k \neq 0$ 且 $\alpha$ 为无理数时，$w =
\eu^{2\iu\pi k\alpha} \neq 1$：
\[
\Bigl|\sum_{n=1}^{N}w^n\Bigr|
= \Bigl|\frac{w(w^N - 1)}{w - 1}\Bigr|
\leq \frac{2}{\abs{1 - w}} ,
\]
与 $N$ 无关的界；除以 $N$ 即得外尔
判别准则，再由问题 17 得结论：$(\{n\alpha\})$ 是
等分布的。

\textbf{19.} 每个子区间都得到渐近比例，
等于其长度，特别地含有无穷多个点：
$(\{n\alpha\})$ 稠密。等分布则更强：一个
序列可以稠密，却几乎把全部时间都花在某个角落
（稠密性说明序列 \emph{去往何处}，
等分布则说明 \emph{频率如何}）。

\textbf{20.} $2^n$ 首位数字为 $1$，当且仅当 $10^m \leq 2^n <
2\cdot10^m$ 对某个 $m$ 成立，即 $\{n\log_{10}2\} \in
\intco{0}{\log_{10}2}$。关于无理性：若 $\log_{10}2 =
\frac pq$，则 $2^q = 10^p = 2^p5^p$，此式对 $p
\geq 1$ 时由唯一分解定理不成立。外尔定理（问题 18
取 $\alpha = \log_{10}2$；半开区间被
夹在长度相近的闭区间之间）给出比例
$\log_{10}2 \approx 0.301$，即 $2^n$ 的首位数字服从
本福特定律。

\textbf{21.} $z$ 为 $C^1$ 类函数，故其傅里叶级数
正规收敛，和为 $z$（\cref{thm:b2:fourier:parseval}~(1)），
且 $c_n(z') = \iu n\,c_n$。因 $\abs{z'} = \frac{L}{2\pi}$
为常数，
\[
\int_0^{2\pi}\abs{z'}^2\dd t
= 2\pi\Bigl(\frac{L}{2\pi}\Bigr)^{\!2}
= \frac{L^2}{2\pi} ,
\]
再将帕塞瓦尔等式用于连续的 $z'$，得
$\frac{1}{2\pi}\int_0^{2\pi}\abs{z'}^2 =
\sum_n\abs{\iu nc_n}^2$，即 $\frac{L^2}{2\pi} = 2\pi\sum_n
n^2\abs{c_n}^2$.

\textbf{22.} 极化形式的帕塞瓦尔等式 $\frac{1}{2\pi}\int\conj
fg = \sum\conj{c_n(f)}c_n(g)$ 可由帕塞瓦尔等式用于
$f + g$ 与 $f + \iu g$（极化恒等式）推出，二者均连续。
取 $f = z$，$g = z'$：
\[
A = \frac12\,\Im\int_0^{2\pi}\conj z\,z'
= \pi\,\Im\sum_n\conj{c_n}(\iu n c_n)
= \pi\sum_n n\abs{c_n}^2 .
\]

\textbf{23.} 综合第 21--22 题：
\[
L^2 - 4\pi A
= 4\pi^2\sum_n n^2\abs{c_n}^2 -
4\pi^2\sum_n n\abs{c_n}^2
= 4\pi^2\sum_{n\in\Z}(n^2 - n)\abs{c_n}^2 \geq 0 ,
\]
因 $n^2 - n = n(n-1) \geq 0$ 对每个整数成立。等号
成立时须 $c_n = 0$（对所有 $n \notin \{0, 1\}$）：$z(t) = c_0 +
c_1\eu^{\iu t}$，是以 $c_0$ 为圆心、
半径为 $\abs{c_1} = \frac{L}{2\pi}$ 的圆（速率为常数）。赫尔维茨
的证明给出等周不等式：$A \leq \frac{L^2}{4\pi}$，
等号仅在圆时成立。

\textbf{24.} 半径为 $R$ 的圆：$L = 2\pi R$，$A = \pi
R^2$：$L^2 = 4\pi^2R^2 = 4\pi A$，取等号。边长为
$a$：$L^2 = 16a^2 > 4\pi a^2 = 4\pi A$（因 $16 > 4\pi \approx
12.57$）。对负整数 $n$，$n^2 - n = n(n - 1)$ 是
两负整数之积，为正数 --- 因此反向绕行的
模式的面积代价加倍。匀速性在第 21 题中引入，
将 $\int\abs{z'}^2$ 化为 $\frac{L^2}{2\pi}$；对
非匀速情形，由柯西–施瓦茨不等式得 $\int\abs{z'}^2 \geq
\frac{(\int\abs{z'})^2}{2\pi} = \frac{L^2}{2\pi}$，故
不等式依然成立，圆仍是唯一的
取等情形。

\textbf{25.} (i) 一切均源于 $F_N \geq 0$（具单位
均值与集中性）；$D_N$ 具单位均值与振荡集中，
但不具正性，吉布斯现象即其代价。
(ii) 傅里叶级数的切萨罗可和性蕴含其
阿贝尔可和性，且和相同（弗罗贝尼乌斯，已在
幂级数一章的周末问题中证明） --- 泊松核
途径（\cref{exo:b2:fourier:12}）恰是阿贝尔方法。
(iii) 外尔定理仅需一致逼近
（第 6 题）；等周不等式则需要帕塞瓦尔等式
本身（第 8、21--22 题）。(iv) 本科三年级分册证明
完备性：指数函数构成 $L^2$ 的希尔伯特基，
帕塞瓦尔等式化为希尔伯特空间的等距映射，费耶尔
定理则表述为：此等距映射可以
由正平均来计算。
\end{solution}
